% GENERATED FILE. Edit canonical lessons, then run npm run generate:books.
% canonical-source-hash: fnv1a64:79f17b6a500fd4df
% canonical-lessons: SA-C402-dharmasala, SA-C402-parnasala, SA-C402-nisidati, SA-C402-samvisati, SA-C402-padyam

\chapter{Rest Houses}
\label{ch:sa-rest-houses}
\hlchaptermodality{402}

\begin{chapteropening}
\textbf{By the end of this chapter:} \emph{I can find a place to rest and sleep on the way.}

Complete the last lesson of chapter 402: I can find a place to rest and sleep on the way.
\end{chapteropening}

\section[\texorpdfstring{\sk{धर्मशाला} (dharmaśālā)}{धर्मशाला (dharmaśālā)}]{\sk{धर्मशाला} (dharmaśālā) — a pilgrims' rest house}

\label{lesson:SA-C402-dharmasala}



\begin{quote}
Before the new one: say the Sanskrit for a lay follower, then the Sanskrit for a memorial shrine.
\end{quote}

\subsection*{You'll want to know: \sk{धर्मशाला}}
\textbf{\sk{धर्मशाला}} — \emph{dharmaśālā} — \textquotedblleft{}a pilgrims' rest house\textquotedblright{}.

\begin{quote}
\sk{यात्रिकाः} \sk{धर्मशालायां} \sk{रात्रिं} \sk{वसन्ति।}
\end{quote}

\begin{quote}
\emph{yātrikāḥ dharmaśālāyāṃ rātriṃ vasanti.}
\end{quote}

\begin{quote}
The pilgrims stay the night in the rest house.
\end{quote}

A feminine noun.

\begin{cousinweb}
\textbf{\sk{धर्मशाला}} (\emph{dharmaśālā}) is built from \emph{dharma}, \textquotedblleft{}right conduct, charity\textquotedblright{}, and \emph{śālā}, \textquotedblleft{}a hall, a house\textquotedblright{}; MW gives the compound as a charitable asylum, especially a religious one (MW).
\end{cousinweb}

\begin{grammarlens}[title={What you've built}]
One more word for getting around.
\end{grammarlens}

\subsection*{Your turn}
\begin{itemize}
\raggedright
  \item \emph{Say it:} \emph{dharmaśālā}
  \item \emph{Say it:} \emph{dharmaśālā}, once more
  \item \emph{Say it:} \emph{yātrikāḥ dharmaśālāyāṃ rātriṃ vasanti.} — The pilgrims stay the night in the rest house.
  \item \emph{Recall:} say the Sanskrit for a lay follower, then the Sanskrit for a memorial shrine, then say \emph{dharmaśālā} again
\end{itemize}

\begin{tcolorbox}[breakable,colback=teal!4,colframe=teal!35!black,title={Before you move on}]
What does \sk{धर्मशाला} mean? (\textquotedblleft{}A pilgrims' rest house\textquotedblright{}.) Say it once more.
\end{tcolorbox}
\section[\texorpdfstring{\sk{पर्णशाला} (parṇaśālā)}{पर्णशाला (parṇaśālā)}]{\sk{पर्णशाला} (parṇaśālā) — a hut of leaves}

\label{lesson:SA-C402-parnasala}



\begin{quote}
Before the new one: say the Sanskrit for a memorial shrine, then the Sanskrit for a pilgrims' rest house.
\end{quote}

\subsection*{You'll want to know: \sk{पर्णशाला}}
\textbf{\sk{पर्णशाला}} — \emph{parṇaśālā} — \textquotedblleft{}a hut of leaves\textquotedblright{}.

\begin{quote}
\sk{परिव्राजकः} \sk{वने} \sk{पर्णशालायां} \sk{वसति।}
\end{quote}

\begin{quote}
\emph{parivrājakaḥ vane parṇaśālāyāṃ vasati.}
\end{quote}

\begin{quote}
The wandering ascetic lives in a hut of leaves in the forest.
\end{quote}

A feminine noun.

\begin{cousinweb}
\textbf{\sk{पर्णशाला}} (\emph{parṇaśālā}) is built from \emph{parṇa}, \textquotedblleft{}a leaf\textquotedblright{}, and \emph{śālā}, \textquotedblleft{}a hall, a house\textquotedblright{}; MW gives the compound as a hut of leaves, an arbour (MW).
\end{cousinweb}

\begin{grammarlens}[title={What you've built}]
One more word for getting around.
\end{grammarlens}

\subsection*{Your turn}
\begin{itemize}
\raggedright
  \item \emph{Say it:} \emph{parṇaśālā}
  \item \emph{Say it:} \emph{parṇaśālā}, once more
  \item \emph{Say it:} \emph{parivrājakaḥ vane parṇaśālāyāṃ vasati.} — The wandering ascetic lives in a hut of leaves in the forest.
  \item \emph{Recall:} say the Sanskrit for a memorial shrine, then the Sanskrit for a pilgrims' rest house, then say \emph{parṇaśālā} again
\end{itemize}

\begin{tcolorbox}[breakable,colback=teal!4,colframe=teal!35!black,title={Before you move on}]
What does \sk{पर्णशाला} mean? (\textquotedblleft{}A hut of leaves\textquotedblright{}.) Say it once more.
\end{tcolorbox}
\section[\texorpdfstring{\sk{निषीदति} (niṣīdati)}{निषीदति (niṣīdati)}]{\sk{निषीदति} (niṣīdati) — sits down to rest}

\label{lesson:SA-C402-nisidati}



\begin{quote}
Before the new one: say the Sanskrit for a pilgrims' rest house, then the Sanskrit for a hut of leaves.
\end{quote}

\subsection*{You'll want to know: \sk{निषीदति}}
\textbf{\sk{निषीदति}} — \emph{niṣīdati} — \textquotedblleft{}sits down to rest\textquotedblright{}.

\begin{quote}
\sk{श्रान्तः} \sk{पथिकः} \sk{वृक्षस्य} \sk{छायायां} \sk{निषीदति।}
\end{quote}

\begin{quote}
\emph{śrāntaḥ pathikaḥ vṛkṣasya chāyāyāṃ niṣīdati.}
\end{quote}

\begin{quote}
The tired wayfarer sits down to rest in the shade of a tree.
\end{quote}

Past (he …): \sk{सः} \sk{छायायां} \sk{निषण्णवान्}, \emph{saḥ chāyāyāṃ niṣaṇṇavān}.

\begin{cousinweb}
\textbf{\sk{निषीदति}} (\emph{niṣīdati}) is built from \emph{ni}, \textquotedblleft{}down\textquotedblright{}, and the root √sad, \textquotedblleft{}to sit\textquotedblright{}; MW gives \emph{ni-ṣad} as to sit down, rest or settle, with the participle \emph{niṣaṇṇa}, \textquotedblleft{}seated, resting\textquotedblright{} (MW).
\end{cousinweb}

\begin{grammarlens}[title={What you've built}]
One more word for getting around.
\end{grammarlens}

\subsection*{Your turn}
\begin{itemize}
\raggedright
  \item \emph{Say it:} \emph{niṣīdati}
  \item \emph{Say it:} \emph{niṣīdati}, once more
  \item \emph{Say it:} \emph{śrāntaḥ pathikaḥ vṛkṣasya chāyāyāṃ niṣīdati.} — The tired wayfarer sits down to rest in the shade of a tree.
  \item \emph{Recall:} say the Sanskrit for a pilgrims' rest house, then the Sanskrit for a hut of leaves, then say \emph{niṣīdati} again
\end{itemize}

\begin{tcolorbox}[breakable,colback=teal!4,colframe=teal!35!black,title={Before you move on}]
What does \sk{निषीदति} mean? (\textquotedblleft{}Sits down to rest\textquotedblright{}.) Say it once more.
\end{tcolorbox}
\section[\texorpdfstring{\sk{संविशति} (saṃviśati)}{संविशति (saṃviśati)}]{\sk{संविशति} (saṃviśati) — lies down to sleep}

\label{lesson:SA-C402-samvisati}



\begin{quote}
Before the new one: say the Sanskrit for a hut of leaves, then the Sanskrit for sits down to rest.
\end{quote}

\subsection*{You'll want to know: \sk{संविशति}}
\textbf{\sk{संविशति}} — \emph{saṃviśati} — \textquotedblleft{}lies down to sleep\textquotedblright{}.

\begin{quote}
\sk{रात्रौ} \sk{यात्रिकाः} \sk{धर्मशालायां} \sk{संविशन्ति।}
\end{quote}

\begin{quote}
\emph{rātrau yātrikāḥ dharmaśālāyāṃ saṃviśanti.}
\end{quote}

\begin{quote}
At night the pilgrims lie down to sleep in the rest house.
\end{quote}

Past (he …): \sk{सः} \sk{धर्मशालायां} \sk{संविष्टवान्}, \emph{saḥ dharmaśālāyāṃ saṃviṣṭavān}.

\begin{cousinweb}
\textbf{\sk{संविशति}} (\emph{saṃviśati}) is built from \emph{sam}, \textquotedblleft{}together\textquotedblright{}, and the root √viś, \textquotedblleft{}to enter, to settle\textquotedblright{}; MW gives \emph{saṃ-viś} as to lie down or go to rest, and its participle \emph{saṃviṣṭa} as one who has lain down, sleeping (MW).
\end{cousinweb}

\begin{grammarlens}[title={What you've built}]
One more word for getting around.
\end{grammarlens}

\subsection*{Your turn}
\begin{itemize}
\raggedright
  \item \emph{Say it:} \emph{saṃviśati}
  \item \emph{Say it:} \emph{saṃviśati}, once more
  \item \emph{Say it:} \emph{rātrau yātrikāḥ dharmaśālāyāṃ saṃviśanti.} — At night the pilgrims lie down to sleep in the rest house.
  \item \emph{Recall:} say the Sanskrit for a hut of leaves, then the Sanskrit for sits down to rest, then say \emph{saṃviśati} again
\end{itemize}

\begin{tcolorbox}[breakable,colback=teal!4,colframe=teal!35!black,title={Before you move on}]
What does \sk{संविशति} mean? (\textquotedblleft{}Lies down to sleep\textquotedblright{}.) Say it once more.
\end{tcolorbox}
\section[\texorpdfstring{\sk{पाद्यम्} (pādyam)}{पाद्यम् (pādyam)}]{\sk{पाद्यम्} (pādyam) — water for washing the feet}

\label{lesson:SA-C402-padyam}



\begin{quote}
Before the new one: say the Sanskrit for sits down to rest, then the Sanskrit for lies down to sleep.
\end{quote}

\subsection*{You'll want to know: \sk{पाद्यम्}}
\textbf{\sk{पाद्यम्}} — \emph{pādyam} — \textquotedblleft{}water for washing the feet\textquotedblright{}.

\begin{quote}
\sk{स्वामी} \sk{अतिथये} \sk{पाद्यं} \sk{ददाति।}
\end{quote}

\begin{quote}
\emph{svāmī atithaye pādyaṃ dadāti.}
\end{quote}

\begin{quote}
The host gives the guest water to wash his feet.
\end{quote}

A neuter noun.

\begin{cousinweb}
\textbf{\sk{पाद्यम्}} (\emph{pādyam}) is built on \emph{pāda}, \textquotedblleft{}a foot\textquotedblright{}, with the ending \emph{-ya}; MW gives \emph{pādya} as belonging to the foot and, as a noun, water for washing the feet, offered to a guest (MW).
\end{cousinweb}

\begin{grammarlens}[title={What you've built}]
One more word for getting around.
\end{grammarlens}

\subsection*{Your turn}
\begin{itemize}
\raggedright
  \item \emph{Say it:} \emph{pādyam}
  \item \emph{Say it:} \emph{pādyam}, once more
  \item \emph{Say it:} \emph{svāmī atithaye pādyaṃ dadāti.} — The host gives the guest water to wash his feet.
  \item \emph{Recall:} say the Sanskrit for sits down to rest, then the Sanskrit for lies down to sleep, then say \emph{pādyam} again
\end{itemize}

\begin{tcolorbox}[breakable,colback=teal!4,colframe=teal!35!black,title={Before you move on}]
What does \sk{पाद्यम्} mean? (\textquotedblleft{}Water for washing the feet\textquotedblright{}.) Say it once more.
\end{tcolorbox}
